% from ega2-1-fr.tex
\oldpage[II]{153}
On notera que ces hypothèses entraînent
\begin{equation}
\label{II.8.1.1.3-fr}
  \mathcal{I}_1^n\subset\mathcal{I}_n.
\tag{8.1.1.3}
\end{equation}

Posons
\begin{equation}
\label{II.8.1.1.4-fr}
  \mathcal{S}=\bigoplus_{n\geq0}\mathcal{I}_n.
\tag{8.1.1.4}
\end{equation}

Il résulte de (\hyperref[II.8.1.1.1-fr]{8.1.1.1}) et
(\hyperref[II.8.1.1.2-fr]{8.1.1.2}) que $\mathcal{S}$ est une
$\mathcal{O}_Y$-Algèbre graduée quasi-cohérente, et définit donc un $Y$-schéma
$X=\operatorname{Proj}(\mathcal{S})$. Si $\mathcal{J}$ est un Idéal
\emph{inversible} de $\mathcal{O}_Y$,
$\mathcal{I}_n\otimes_{\mathcal{O}_Y}\mathcal{J}^{\otimes n}$ s'identifie
canoniquement à $\mathcal{I}_n\mathcal{J}^n$. Si on remplace donc les
$\mathcal{I}_n$ par les $\mathcal{I}_n\mathcal{J}^n$, ce qui remplace
$\mathcal{S}$ par une $\mathcal{O}_Y$-Algèbre quasi-cohérente
$\mathcal{S}_{(\mathcal{J})}$,
$X_{(\mathcal{J})}=\operatorname{Proj}(\mathcal{S}_{(\mathcal{J})})$ est
canoniquement isomorphe à $X$ (\hyperref[II.3.1.8-fr]{3.1.8}).
\end{env}

\begin{env}[8.1.2]
\label{II.8.1.2-fr}
Supposons $Y$ \emph{localement intègre}, de sorte que le faisceau
$\mathcal{R}(Y)$ des fonctions rationnelles est une $\mathcal{O}_Y$-Algèbre
quasi-cohérente (\hyperref[I.7.3.7-fr]{I, 7.3.7}). Nous dirons qu'un
sous-$\mathcal{O}_Y$-Module $\mathcal{I}$ de $\mathcal{R}(Y)$ est un
\emph{Idéal fractionnaire} de $\mathcal{R}(Y)$ s'il est de \emph{type fini}
(\hyperref[0.5.2.1-fr]{0, 5.2.1}). Supposons donné, pour tout $n\geq0$, un
Idéal fractionnaire quasi-cohérent $\mathcal{I}_n$ de $\mathcal{R}(Y)$, tel
que $\mathcal{I}_0=\mathcal{O}_Y$, et que la condition
(\hyperref[II.8.1.1.2-fr]{8.1.1.2}) soit vérifiée (mais non nécessairement la
seconde condition (\hyperref[II.8.1.1.1-fr]{8.1.1.1}))~; on peut encore alors
définir par la formule (\hyperref[II.8.1.1.4-fr]{8.1.1.4}) une
$\mathcal{O}_Y$-Algèbre graduée quasi-cohérente, et le $Y$-schéma correspondant
$X=\operatorname{Proj}(\mathcal{S})$~; on aura encore un isomorphisme
canonique de $X$ sur $X_{(\mathcal{J})}$ pour tout Idéal fractionnaire
inversible $\mathcal{J}$ de $\mathcal{R}(Y)$.
\end{env}

\begin{definition}[8.1.3]
\label{II.8.1.3-fr}
Soit $Y$ un préschéma (resp. un préschéma localement intègre), et soit
$\mathcal{I}$ un Idéal quasi-cohérent de $\mathcal{O}_Y$ (resp. un Idéal
fractionnaire quasi-cohérent de $\mathcal{R}(Y)$). On dit que le $Y$-schéma
$X=\operatorname{Proj}(\bigoplus_{n\geq0}\mathcal{I}^n)$ est obtenu en
\emph{faisant éclater l'Idéal $\mathcal{I}$}, ou est le \emph{préschéma éclaté
de $Y$ relativement à $\mathcal{I}$}. Lorsque $\mathcal{I}$ est un Idéal
quasi-cohérent de $\mathcal{O}_Y$ et $Y'$ le sous-préschéma fermé de $Y$ défini
par $\mathcal{I}$, on dit aussi que $X$ est le $Y$-schéma obtenu en faisant
éclater $Y'$.
\end{definition}

Par définition, $\mathcal{S}=\bigoplus_{n\geq0}\mathcal{I}^n$ est alors
engendrée par $\mathcal{S}_1=\mathcal{I}$~; si $\mathcal{I}$ est un
$\mathcal{O}_Y$-Module de \emph{type fini}, $X$ est donc \emph{projectif} sur
$Y$ (\hyperref[II.5.5.2-fr]{5.5.2}). Sans hypothèse sur $\mathcal{I}$, le
$\mathcal{O}_X$-Module $\mathcal{O}_X(1)$ est \emph{inversible}
(\hyperref[II.3.2.5-fr]{3.2.5}) et \emph{très ample} en vertu de
(\hyperref[II.4.4.3-fr]{4.4.3}) pour le morphisme structural $X\to Y$.

On notera que si $f:X\to Y$ est le morphisme structural, la restriction de
$f$ à $f^{-1}(Y-Y')$ est un \emph{isomorphisme} sur $Y-Y'$ lorsque
$\mathcal{I}$ est un Idéal de $\mathcal{O}_Y$ et $Y'$ le sous-préschéma fermé
qu'il définit~: en effet, la question étant locale sur $Y$, il suffit de
supposer $\mathcal{I}=\mathcal{O}_Y$, et notre assertion résulte de
(\hyperref[II.3.1.7-fr]{3.1.7}).

Lorsqu'on remplace $\mathcal{I}$ par $\mathcal{I}^d$ ($d>0$), le $Y$-schéma
éclaté $X$ est remplacé par un $Y$-schéma canoniquement isomorphe $X'$
(\hyperref[II.8.1.1-fr]{8.1.1})~; de même, pour tout Idéal (resp. Idéal
fractionnaire) inversible $\mathcal{J}$, le préschéma éclaté
$X_{(\mathcal{J})}$ relativement à l'Idéal $\mathcal{I}\mathcal{J}$ est
canoniquement isomorphe à $X$ (\hyperref[II.8.1.1-fr]{8.1.1}).

En particulier, lorsque $\mathcal{I}$ est un Idéal (resp. un Idéal
fractionnaire) \emph{inversible}, le $Y$-schéma obtenu en faisant éclater
$\mathcal{I}$ est \emph{isomorphe} à $Y$ (\hyperref[II.3.1.7-fr]{3.1.7}).

\begin{proposition}[8.1.4]
\label{II.8.1.4-fr}
Soit $Y$ un préschéma intègre.
\begin{enumerate}
\item[\rm(i)] Pour toute suite $(\mathcal{I}_n)$ d'Idéaux fractionnaires
quasi-cohérents de $\mathcal{R}(Y)$ vérifiant
(\hyperref[II.8.1.1.2-fr]{8.1.1.2})

